\section{AHSME 1985}
        \begin{enumerate}
        	\item (\textit{AHSME 1985}) If $2x+1=8$, then $4x+1=$


 $\mathrm{(A)\ } 15 \qquad \mathrm{(B) \ }16 \qquad \mathrm{(C) \  } 17 \qquad \mathrm{(D) \  } 18 \qquad \mathrm{(E) \  }19$


 
 \textbf{Answer: }A



\textbf{Solution 1}

We have 
 \begin{align*}2x+1 = 8 &\iff 2x = 7 \\ &\iff x = \frac{7}{2},\end{align*}
 so 
 \begin{align*}4x+1 &= 4\left(\frac{7}{2}\right)+1 \\ &= 2(7)+1 \\ &= \boxed{\text{(A)} \ 15}.\end{align*}


 



\textbf{Solution 2}

From $2x = 7$ (as above), we can directly compute 
 \begin{align*}4x &= 2(2x) \\ &= 2(7) \\ &= 14,\end{align*}
 so $4x+1 = 14+1 = \boxed{\text{(A)} \ 15}$.

 



\textbf{Solution 3}

Multiply both sides of the equation by $2$ to get $4x + 2 = 16$.  Then subtract $1$ from both sides to get $4x + 1 = \boxed{\text{(A)} \ 15}$

 -j314andrews

 


	\item (\textit{AHSME 1985}) In an arcade game, the "monster" is the shaded sector of a circle of radius $1$ cm, as shown in the figure. The missing piece (the mouth) has central angle $ 60\degree$. What is the perimeter of the monster in cm?


 
                    \begin{figure}[H]
                        centering
                        \includegraphics[width=0.5\textwidth]{Problem/AHSME_I_1985_2_Problem_Figure_1.png}
                    \end{figure}
                    


 $\mathrm{(A)\ } \pi+2 \qquad \mathrm{(B) \ }2\pi \qquad \mathrm{(C) \  } \frac{5}{3}\pi \qquad \mathrm{(D) \  } \frac{5}{6}\pi+2 \qquad \mathrm{(E) \  }\frac{5}{3}\pi+2$


 
 \textbf{Answer: }E



\textbf{Solution}

The two straight sides of the "monster" are radii of length $1$ cm, so their total length is $1+1 = 2$ cm. The circular arc forms $\frac{360^{\circ}-60^{\circ}}{360^{\circ}} = \frac{5}{6}$ of the entire circle, and a circle with radius $1$ cm has circumference $2\pi(1) = 2\pi$ cm, so the length of the arc is $\left(\frac{5}{6}\right)(2\pi) = \frac{5}{3}\pi$ cm. Hence the total perimeter in cm is $\boxed{\text{(E)} \ \frac{5}{3}\pi+2}$.

 


	\item (\textit{AHSME 1985}) In right $\triangle ABC$ with legs $5$ and $12$, arcs of circles are drawn, one with center $A$ and radius $12$, the other with center $B$ and radius $5$. They intersect the hypotenuse in $M$ and $N$. Then $MN$ has length 


 
                    \begin{figure}[H]
                        centering
                        \includegraphics[width=0.5\textwidth]{Problem/AHSME_I_1985_3_Problem_Figure_1.png}
                    \end{figure}
                    


 $\mathrm{(A)\ } 2 \qquad \mathrm{(B) \ }\frac{13}{5} \qquad \mathrm{(C) \  } 3 \qquad \mathrm{(D) \  } 4 \qquad \mathrm{(E) \  }\frac{24}{5}$


 
 \textbf{Answer: }D



\textbf{Solution}

Firstly, the Pythagorean theorem gives 
 \begin{align*}AB &=\sqrt{AC^2+BC^2} \\ &= \sqrt{12^2+5^2} \\ & =\sqrt{144+25} \\ &=\sqrt{169} \\ &= 13.\end{align*}
 Also, $AM = AC = 12$ and $BN = BC = 5$ since they are both radii of the respective circles. Thus $MB = AB-AM = 13-12 = 1$, and so $MN = BN-BM = 5-1 = \boxed{\text{(D)} \ 4}$.

 


	\item (\textit{AHSME 1985}) A large bag of coins contains pennies, dimes and quarters. There are twice as many dimes as pennies and three times as many quarters as dimes. An amount of money which could be in the bag is


 $\mathrm{(A)\ } \$306 \qquad \mathrm{(B) \ }  \$333 \qquad \mathrm{(C)\ } \$342 \qquad \mathrm{(D) \  }  \$348 \qquad \mathrm{(E) \  }  \$360$


 
 \textbf{Answer: }C



\textbf{Solution}

If there are $x$ pennies in the bag, then there are $2x$ dimes and $3(2x) = 6x$ quarters. Since pennies are $\$0.01$, dimes are $\$0.10$, and quarters are $\$0.25$, the total amount of money in the bag is 
 \[\$ \left(0.01x+(0.10)(2x)+(0.25)(6x)\right) = \$1.71x.\] Therefore, the possible amounts of money are precisely the integer multiples of $\$1.71$. 

 Since the answer choices are all integer numbers of dollars, we multiply by $100$ to deduce that the answer must be an integer multiple of $\$171$. The only such multiple among the answer choices is $\$(2 \cdot 171) = \boxed{\text{(C)} \ \$342}$.

 


	\item (\textit{AHSME 1985}) Which terms must be removed from the sum


 
 \[\frac{1}{2}+\frac{1}{4}+\frac{1}{6}+\frac{1}{8}+\frac{1}{10}+\frac{1}{12}\]


 if the sum of the remaining terms is to equal $1$?


 $\mathrm{(A)\ } \frac{1}{4}\text{ and }\frac{1}{8} \qquad \mathrm{(B) \ }\frac{1}{4}\text{ and }\frac{1}{12} \qquad \mathrm{(C) \  } \frac{1}{8}\text{ and }\frac{1}{12} \qquad \mathrm{(D) \  } \frac{1}{6}\text{ and }\frac{1}{10} \qquad \mathrm{(E) \  }\frac{1}{8}\text{ and }\frac{1}{10}$


 
 \textbf{Answer: }E



\textbf{Solution}

We compute the entire sum as 
 \begin{align*}\frac{1}{2}+\frac{1}{4}+\frac{1}{6}+\frac{1}{8}+\frac{1}{10}+\frac{1}{12} &= \frac{60+30+20+15+12+10}{120} \\ &= \frac{147}{120} \\ &= \frac{49}{40},\end{align*}
so the two terms to be removed must sum to $\frac{49}{40}-1 = \frac{9}{40}$. That is, $\frac{9}{40} = \frac{1}{x}+\frac{1}{y}=\frac{x+y}{xy}$ for some $x,y \in \{2,4,6,8,10,12\}$.

 Since $9$ and $40$ are coprime, we deduce that $xy$ is a multiple of $40$, so $x$ or $y$ must be a multiple of $5$. It follows that one of them must be $10$, and we obtain 
 \begin{align*}\frac{9}{40}-\frac{1}{10} &= \frac{9-4}{40} \\ &= \frac{5}{40} \\ &= \frac{1}{8},\end{align*}
 so the two terms we must remove are $\boxed{\text{(E)} \ \frac{1}{8}\text{ and }\frac{1}{10}}$.

 


	\item (\textit{AHSME 1985}) One student in a class of boys and girls is chosen to represent the class. Each student is equally likely to be chosen and the probability that a boy is chosen is $\frac{2}{3}$ of the probability that a girl is chosen. The ratio of the number of boys to the total number of boys and girls is


 $\mathrm{(A)\ } \frac{1}{3} \qquad \mathrm{(B) \ }\frac{2}{5} \qquad \mathrm{(C) \  } \frac{1}{2} \qquad \mathrm{(D) \  } \frac{3}{5} \qquad \mathrm{(E) \  }\frac{2}{3}$


 
 \textbf{Answer: }B



\textbf{Solution}

Let the probability that a boy is chosen be $p$, so the probability that a girl is chosen is $1-p$ (as the probabilities must sum to $1$). Therefore 
 \begin{align*}\frac{2}{3}(1-p) = p &\iff 2(1-p) = 3p \\ &\iff 2-2p = 3p \\&\iff 5p = 2 \\&\iff p = \frac{2}{5}.\end{align*}


 Now simply notice that the probability a boy is chosen is precisely the proportion of the number of boys out of the total number of students (boys and girls), so the answer is precisely $\boxed{\text{(B)} \ \frac{2}{5}}$.

 


	\item (\textit{AHSME 1985}) In some computer languages (such as APL), when there are no parentheses in an algebraic expression, the operations are grouped from right to left. Thus, $a \times b - c$ in such languages means the same as $a(b-c)$ in ordinary algebraic notation. If $a\div b-c+d$ is evaluated in such a language, the result in ordinary algebraic notation would be


 $\mathrm{(A)\ } \frac{a}{b}-c+d \qquad \mathrm{(B) \ }\frac{a}{b}-c-d \qquad \mathrm{(C) \  } \frac{d+c-b}{a} \qquad \mathrm{(D) \  } \frac{a}{b-c+d} \qquad \mathrm{(E) \  }\frac{a}{b-c-d}$


 
 \textbf{Answer: }E



\textbf{Solution}

The rightmost part of the expression is $c+d$, so $b-c+d$ would be grouped as $b-(c+d)$, and thus the whole expression would be grouped as $a\div (b-(c+d)) = \boxed{\text{(E)} \ \frac{a}{b-c-d}}$.

 


	\item (\textit{AHSME 1985}) Let $a,a',b,b'$ be real numbers with $a$ and $a'$ nonzero. The solution to $ax+b=0$ is less than the solution to $a'x+b'=0$ if and only if 


 $\mathrm{(A)\ } a'b < ab' \qquad \mathrm{(B) \ }ab' < a'b \qquad \mathrm{(C) \  } ab < a'b' \qquad \mathrm{(D) \  } \frac{b}{a} < \frac{b'}{a'} \qquad \mathrm{(E) \  }\frac{b'}{a'} < \frac{b}{a}$


 
 \textbf{Answer: }E



\textbf{Solution}

The solution to $ax+b=0$ is $x = \frac{-b}{a}$, while that to $a'x+b'=0$ is $x = \frac{-b'}{a'}$. The first solution is less than the second precisely if 
 \[\frac{-b}{a} < \frac{-b'}{a'},\] and multiplying this inequality by $-1$ reversees the inequality sign, yielding $\boxed{\text{(E)} \ \frac{b'}{a'} < \frac{b}{a}}$.

 


	\item (\textit{AHSME 1985}) The odd positive integers $1, 3, 5, 7, \ldots$, are arranged into five columns continuing with the pattern shown on the right. Counting from the left, the column in which $1985$ appears is the


 
                    \begin{figure}[H]
                        centering
                        \includegraphics[width=0.5\textwidth]{Problem/AHSME_I_1985_9_Problem_Figure_1.png}
                    \end{figure}
                    


 $\mathrm{(A)\ } \text{first} \qquad \mathrm{(B) \ }\text{second} \qquad \mathrm{(C) \  } \text{third} \qquad \mathrm{(D) \  } \text{fourth} \qquad \mathrm{(E) \  }\text{fifth}$


 
 \textbf{Answer: }B



\textbf{Solution}

Considering each integer modulo $16$ gives the following pattern:
                    \begin{figure}[H]
                        centering
                        \includegraphics[width=0.5\textwidth]{Solution/AHSME_I_1985_9_Solution_Figure_1.png}
                    \end{figure}
                    

 We therefore observe that all numbers congruent to $1 \pmod{16}$ will appear in the second column, and since $1985 \equiv 1 \pmod{16}$, the answer is $\boxed{\text{(B)} \ \text{second}}$.

 


	\item (\textit{AHSME 1985}) An arbitrary circle can intersect the graph of $y = \sin x$ in


 $\mathrm{(A)\  } \text{at most }2\text{ points} \qquad \mathrm{(B)\  }\text{at most }4\text{ points} \qquad \mathrm{(C) \  } \text{at most }6\text{ points} \qquad \mathrm{(D) \ } \text{at most }8\text{ points}$
$\mathrm{(E) \  }\text{more than }16\text{ points}$


 
 \textbf{Answer: }E



\textbf{Solution}

Consider a circle whose center lies on the positive $y$-axis and which passes through the origin. As the radius of this circle becomes arbitrarily large, its curvature near the $x$-axis becomes almost flat, and so it can intersect the curve $y = \sin x$ arbitrarily many times (since the $x$-axis itself intersects the curve infinitely many times). Hence, in particular, we can choose a radius sufficiently large that the circle intersects the curve at $\boxed{\text{(E)} \ \text{more than }16 \text{ points}}$.

 


	\item (\textit{AHSME 1985}) How many distinguishable rearrangements of the letters in $CONTEST$ have both the vowels first? (For instance, $OETCNST$ is one such arrangement, but $OTETSNC$ is not.)


 $\mathrm{(A)\ } 60 \qquad \mathrm{(B) \ }120 \qquad \mathrm{(C) \  } 240 \qquad \mathrm{(D) \  } 720 \qquad \mathrm{(E) \  }2520$


 
 \textbf{Answer: }B



\textbf{Solution}

We consider the vowels and consonants separately. There are $2$ vowels ($O$ and $E$), giving $2! = 2$ choices for the first two letters; similarly, there are $5$ consonants ($C$, $N$, $S$, and two $T$s), which would give $5! = 120$ possible choices for letters $3$ to $7$, except that since the two $T$s are indistinguishable, this actually counts each order exactly twice. Therefore the number of possible orderings of the consonants is $\frac{120}{2} = 60$, giving a total of $2 \cdot 60 = \boxed{\text{(B)} \ 120}$ possible rearrangements.

 


	\item (\textit{AHSME 1985}) Let $p$, $q$ and $r$ be distinct prime numbers, where $1$ is not considered a prime. Which of the following is the smallest positive perfect cube having $n = pq^2r^4$ as a divisor?


 $\mathrm{(A)\ } p^8q^8r^8 \qquad \mathrm{(B) }\left(pq^2r^2\right)^3 \qquad \mathrm{(C) } \left(p^2q^2r^2\right)^3 \qquad \mathrm{(D) } \left(pqr^2\right)^3 \qquad \mathrm{(E) \  }4p^3q^3r^3$


 
 \textbf{Answer: }D



\textbf{Solution}

A number of the form $p^aq^br^c$ will be a perfect cube precisely when $a$, $b$, and $c$ are multiples of 3. (Clearly, since we are looking for the \textit{smallest} possible perfect cube, we can assume that it has no other prime factors.) Furthermore, for it to be a multiple of $pq^2r^4$, we must have $a \geq 1$, $b \geq 2$, and $c \geq 4$. The smallest multiple of $3$ that is at least $1$ is $3$, the smallest that is at least $2$ is again $3$, and the smallest that is at least $4$ is $6$. Hence the smallest possible perfect cube with $n$ as a divisor is $p^3q^3r^6 = \boxed{\text{(D)} \left(pqr^2\right)^3}$.

 


	\item (\textit{AHSME 1985}) Pegs are put in a board $1$ unit apart both horizontally and vertically. A rubber band is stretched over $4$ pegs as shown in the figure, forming a quadrilateral. Its area in square units is


 
                    \begin{figure}[H]
                        centering
                        \includegraphics[width=0.5\textwidth]{Problem/AHSME_I_1985_13_Problem_Figure_1.png}
                    \end{figure}
                    


 $\mathrm{(A)\ } 4 \qquad \mathrm{(B) \ }4.5 \qquad \mathrm{(C) \  } 5 \qquad \mathrm{(D) \  } 5.5 \qquad \mathrm{(E) \  }6$


 
 \textbf{Answer: }E



\textbf{Solution 1}


                    \begin{figure}[H]
                        centering
                        \includegraphics[width=0.5\textwidth]{Solution/AHSME_I_1985_13_Solution_Figure_1.png}
                    \end{figure}
                    We draw in the rectangle bounding the given quadrilateral and label the points as shown. The area of rectangle $ABCD$ is $(3)(4) = 12$, while the areas of the triangles $AEH$, $EBF$, $FCG$, and $GDH$ are, respectively, 
 \begin{align*}&\text{area of } \triangle AEH = \frac{1}{2}(1)(2) = 1, \\ &\text{area of } \triangle EBF = \frac{1}{2}(3)(2) = 3, \\ &\text{area of } \triangle FCG = \frac{1}{2}(1)(1) = \frac{1}{2}, \text{ and} \\ &\text{area of } \triangle GDH = \frac{1}{2}(3)(1) = \frac{3}{2}.\end{align*}
Hence the area of the given quadrilateral $EFGH$ is $12-\left(1+3+\frac{1}{2}+\frac{3}{2}\right) = \boxed{\text{(E)} \ 6}$.

 



\textbf{Solution 2}

The number of lattice points (i.e. pegs on the board) strictly inside the quadrilateral is $5$ and the number of lattice points on its boundary is $4$. Therefore, by \href{https://artofproblemsolving.com/wiki/index.php?title=Pick\%27s\_theorem}{Pick's theorem}, its area is $5+\frac{4}{2}-1 = \boxed{\text{(E)} \ 6}$.

 


	\item (\textit{AHSME 1985}) Exactly three of the interior angles of a convex polygon are obtuse. What is the maximum number of sides of such a polygon?


 $\mathrm{(A)\ } 4 \qquad \mathrm{(B) \ }5 \qquad \mathrm{(C) \  } 6 \qquad \mathrm{(D) \  } 7 \qquad \mathrm{(E) \  }8$


 
 \textbf{Answer: }C



\textbf{Solution}

Suppose that such a polygon has $n$ sides. Let the three obtuse angle measures, in degrees, be $o_1$, $o_2$, and $o_3$ and the $(n-3)$ acute angle measures, again in degrees, be $a_1,a_2,a_3, \dotsc a_{n-3}$.

 Since $90 < o_i < 180$ for each $i$, we have 
 \[3(90) = 270 < o_1+o_2+o_3 < 3(180) = 540,\] and similarly, since $0 < a_i < 90$ for each $i$, 
 \[0 < a_1+a_2+a_3+\dotsb+a_{n-3} < 90(n-3) = 90n-270.\]It follows that 
 \[270+0 < o_1+o_2+o_3+a_1+a_2+a_3+\dotsb+a_{n-3} < 540+90n-270,\] and recalling that the sum of the interior angle measures of an $n$-gon is $180(n-2) = 180n-360$, this reduces to $270 < 180n-360 < 90n+270$. Hence 
 \[\frac{540}{180} < n < \frac{270+360}{90} \iff 3 < n < 7,\] so an upper bound is $n \leq 6$, and it is easy to check that this bound can be attained by e.g. a convex hexagon with a right angle, $2$ acute angles, and $3$ obtuse angles, as shown below:

 
                    \begin{figure}[H]
                        centering
                        \includegraphics[width=0.5\textwidth]{Solution/AHSME_I_1985_14_Solution_Figure_1.gif}
                    \end{figure}
                    

 Accordingly, the maximum possible number of sides of such a polygon is $\boxed{\text{(C)} \ 6}$.

 


	\item (\textit{AHSME 1985}) If $a$ and $b$ are positive numbers such that $a^b = b^a$ and $b = 9a$, then the value of $a$ is


 $\mathrm{(A)\ } 9 \qquad \mathrm{(B) \ }\frac{1}{9} \qquad \mathrm{(C) \  } \sqrt[9]{9} \qquad \mathrm{(D) \  } \sqrt[3]{9} \qquad \mathrm{(E) \  }\sqrt[4]{3}$


 
 \textbf{Answer: }E



\textbf{Solution}

Substituting $b = 9a$ into $a^b = b^a$ gives 
 \begin{align*}a^{9a} = \left(9a\right)^a &\iff \left(a^9\right)^a = \left(9a\right)^a \qquad \text{(using the identity } \left(x^y\right)^z = x^{yz}\text{ for } x > 0\text{)} \\ &\iff a^9 = 9a \qquad \text{(taking the } a\text{th root of both sides, as } a > 0\text{)} \\ &\iff a^8 = 9 \qquad \text{(as } a \neq 0\text{)} \\ &\iff a^4 = 3 \qquad \text{(taking square roots and noting that } a^4 \geq 0\text{}) \\ &\iff a = \boxed{\text{(E)} \ \sqrt[4]{3}} \qquad \text{(again as } a > 0\text{)}.\end{align*}


 


	\item (\textit{AHSME 1985}) If $ A = 20 \degree$ and $ B = 25 \degree$, then the value of $\left(1+\tan A\right)\left(1+\tan B\right)$ is


 $\mathrm{(A)\ } \sqrt{3} \qquad \mathrm{(B) \ }2 \qquad \mathrm{(C) \  } 1+\sqrt{2} \qquad \mathrm{(D) \  } 2\left(\tan A+\tan B\right) \qquad \mathrm{(E) \  }\text{none of these}$


 
 \textbf{Answer: }B



\textbf{Solution 1}

Noting that $ 25 \degree + 20 \degree = 45 \degree$, we apply the angle sum formula 
 \[\tan(A+B) = \frac{\tan A+\tan B}{1-\tan A\tan B},\] giving 
 \begin{align*}1 &= \tan 45^{\circ} \\ &= \tan(A+B) \\ &= \frac{\tan A+\tan B}{1-\tan A\tan B},\end{align*}
 so 
 \[\tan A + \tan B = 1-\tan A\tan B.\] Hence 
 \begin{align*}(1+\tan A)(1+\tan B) &= 1+\tan A+\tan B+\tan A\tan B \\ &= 1+\left(1-\tan A\tan B\right)+\tan A\tan B \\ &= \boxed{\text{(B)} \ 2}.\end{align*}


 



\textbf{Solution 2}

Expanding in terms of sines and cosines, we obtain 
 \begin{align*}\left(1+\tan A\right)\left(1+\tan B\right) &= \left(\frac{\cos A}{\cos A}+\frac{\sin A}{\cos A}\right)\left(\frac{\cos B}{\cos B}+\frac{\sin B}{\cos B}\right) \\ &=\frac{\left(\sin A+\cos A\right)\left(\sin B+\cos B\right)}{\cos A\cos B} \\ &= \frac{\sin A\sin B+\cos A\cos B+\sin A\cos B+\sin B\cos A}{\cos A\cos B}.\end{align*}


 Recalling the angle sum identities 
 \begin{align*}&\cos(A-B) = \sin A\sin B+\cos A\cos B \text{ and} \\ &\sin(A+B)=\sin A\cos B+\sin B\cos A,\end{align*}
 this reduces to 
 \[\frac{\cos(A-B)+\sin(A+B)}{\cos A\cos B}.\]

 Now, using the product-to-sum formula 
 \[\cos A\cos B = \frac{1}{2}\left(\cos(A-B)+\cos(A+B)\right),\] we can simplify the denominator, yielding 
 \[\left(1+\tan A\right)\left(1+\tan B\right) = \frac{\cos(A-B)+\sin(A+B)}{\frac{1}{2}\left(\cos(A-B)+\cos(A+B)\right)}.\]

 Finally, since $ A+B = 45 \degree$, we have $\sin(A+B) = \cos(A+B)$, so
 \begin{align*}\left(1+\tan A\right)\left(1+\tan B\right) &= \frac{\cos(A-B)+\sin(A+B)}{\frac{1}{2}\left(\cos(A-B)+\sin(A+B)\right)} \\ &= \frac{1}{\left(\frac{1}{2}\right)} \\ &= \boxed{\text{(B)} \ 2}.\end{align*}


 \textit{Remark}: Notice that we only used the fact that $\sin(A+B) = \cos(A+B)$, so we have in fact shown that $\left(1+\tan A\right)\left(1+\tan B\right) = 2$ not just for $ A = 20 \degree$ and $ B = 25 \degree$, but also for all $A,B$ such that $ A+B = 45 \degree + 180 \degree n$ for integers $n$.

 



\textbf{Solution 3}

As in Solution 2, we rewrite the expression as 
 \[\frac{\cos 20^{\circ} \cos 25^{\circ} + \cos 20^{\circ} \sin 25^{\circ} + \sin 20^{\circ} \cos 25^{\circ} + \sin 20^{\circ} \sin 25^{\circ}}{\cos 20^{\circ} \cos 25^{\circ}},\] and hence as 
 \[1 + \frac{\cos 20^{\circ} \sin 25^{\circ} + \sin 20^{\circ} \cos 25^{\circ} + \sin 20^{\circ} \sin 25^{\circ}}{\cos 20^{\circ} \cos 25^{\circ}}.\] Using the angle sum identities 
 \begin{align*}&\cos(A-B) = \sin A\sin B+\cos A\cos B \text{ and} \\ &\sin(A+B)=\sin A\cos B+\sin B\cos A,\end{align*}
 we obtain 
 \begin{align*}&\cos 20^{\circ} \sin 25^{\circ} + \sin 20^{\circ} \cos 25^{\circ} = \sin\left(20^{\circ}+25^{\circ}\right) = \sin 45^{\circ} \text{ and} \\ &\cos 20^{\circ} \cos 25^{\circ} - \sin 20^{\circ} \sin 25^{\circ} = \cos\left(20^{\circ}+25^{\circ}\right) = \cos 45^{\circ}.\end{align*}
 Therefore the expression becomes 
 \begin{align*}1+\frac{\sin 45^{\circ} + \sin 20^{\circ} \sin 25^{\circ}}{\cos 45^{\circ} + \sin 20^{\circ} \sin 25^{\circ}} &= 1+1 \qquad \text{(since } \sin 45^{\circ} = \cos 45^{\circ}\text{)} \\ &= \boxed{\text{(B)} \ 2}.\end{align*}


 



\textbf{Solution 4}

As in Solutions 2 and 3, the expression becomes 
 \[\left(\frac{\cos A+\sin A}{\cos A}\right)\left(\frac{\cos B+\sin B}{\cos B}\right).\]

 Now, using the identity $\cos^2 A + \sin^2 A = 1$ and the double-angle identity $\sin(2A) = 2\sin A\cos A$, we observe that 
 \begin{align*}\left(\cos A + \sin A\right)^2 &= \cos^2 A + \sin^2 A + 2\sin A \cos A \\ &= 1 + 2\sin A \cos A \\ &= 1 + \sin(2A).\end{align*}


 Since $A$ and $B$ are acute, we have $\sin A,\cos A,\sin B,\cos B > 0$, so $\cos A + \sin A > 0$ and $\cos B + \sin B > 0$. Hence, taking the positive square root of both sides in the above identity, the expression becomes 
 \[\left(\frac{\sqrt{1+\sin(2A)}}{\cos A}\right)\left(\frac{\sqrt{1+\sin(2B)}}{\cos B}\right).\] Recalling the further identity $\sin A = \cos\left(90^{\circ}-A\right)$, together with the half-angle identity 
 \[\cos\left(\frac{A}{2}\right) = \sqrt{\frac{1+\cos A}{2}} \qquad \text{for } 0^{\circ} \leq A \leq 180^{\circ},\] we finally obtain 
 \begin{align*}\left(1+\tan A\right)\left(1+\tan B\right) &= 2\left(\frac{\sqrt{\frac{1+\cos\left(90^{\circ}-2A\right)}{2}}}{\cos A}\right)\left(\frac{\sqrt{\frac{1+\cos\left(90^{\circ}-2B\right)}{2}}}{\cos B}\right) \\ &= 2\left(\frac{\cos\left(45^{\circ}-A\right)}{\cos A}\right)\left(\frac{\cos\left(45^{\circ}-B\right)}{\cos B}\right) \\ &= 2\left(\frac{\cos 25^{\circ} \cos 20^{\circ}}{\cos 20^{\circ} \cos 25^{\circ}}\right) = \boxed{\text{(B)} \ 2}.\end{align*}


 


	\item (\textit{AHSME 1985}) Diagonal $DB$ of rectangle $ABCD$ is divided into three segments of length $1$ by parallel lines $L$ and $L'$ that pass through $A$ and $C$ and are perpendicular to $DB$. The area of $ABCD$, rounded to the one decimal place, is 


 
                    \begin{figure}[H]
                        centering
                        \includegraphics[width=0.5\textwidth]{Problem/AHSME_I_1985_17_Problem_Figure_1.png}
                    \end{figure}
                    


 $\mathrm{(A)\ } 4.1 \qquad \mathrm{(B) \ }4.2 \qquad \mathrm{(C) \  } 4.3 \qquad \mathrm{(D) \  } 4.4 \qquad \mathrm{(E) \  }4.5$


 
 \textbf{Answer: }B



\textbf{Solution}

Let $E$ be the point of intersection of $L$ and $\overline{BD}$. Then, because $AE$ is the altitude to the hypotenuse of right triangle $ABD$, triangles $ADE$ and $BAE$ are similar, giving 
 \[\frac{AE}{BE} = \frac{ED}{EA},\] and so 
 \begin{align*}AE &= \sqrt{BE \cdot ED} \\ &= \sqrt{(1+1)(1)} \\ &= \sqrt{2}.\end{align*}
 Thus, taking $BD$ and $AE$ as the base and perpendicular height, respectively, of triangle $ABD$, we may compute its area as $\frac{1}{2}(3)\left(\sqrt{2}\right) = \frac{3\sqrt{2}}{2}$. By symmetry, the area of the entire rectangle $ABCD$ is 
 \[2\left(\frac{3\sqrt{2}}{2}\right) = 3\sqrt{2} \approx (3)(1.4) = \boxed{\text{(B)} \ 4.2}.\]

 


	\item (\textit{AHSME 1985}) Six bags of marbles contain $18$, $19$, $21$, $23$, $25$ and $34$ marbles, respectively. One bag contains chipped marbles only. The other $5$ bags contain no chipped marbles. Jane takes three of the bags and George takes two of the others. Only the bag of chipped marbles remains. If Jane gets twice as many marbles as George, how many chipped marbles are there?


 $\mathrm{(A)\ } 18 \qquad \mathrm{(B) \ }19 \qquad \mathrm{(C) \  } 21 \qquad \mathrm{(D) \  } 23 \qquad \mathrm{(E) \  }25$


 
 \textbf{Answer: }D



\textbf{Solution}

Let George's bags contain a total of $x$ marbles, so Jane's bag contains $2x$ marbles. This means the total number of non-chipped marbles is $3x \equiv 0 \pmod{3}$, while the total number of marbles is $18+19+21+23+25+34 = 140 \equiv 2 \pmod{3}$, so the number of chipped marbles must also be congruent to $2-0 \equiv 2 \pmod{3}$.

 The answer choices are congruent modulo 3 to $0$, $1$, $0$, $2$, and $1$ respectively, so the only possible number of chipped marbles among these is $23$. Indeed, if Jane takes the bags containing $19$, $25$, and $34$ marbles and George takes the remaining bags containing $18$ and $21$ marbles, then Jane will have a total of $19+25+34 = 78$ marbles, which is twice as many as George's $18+21 = 39$ marbles, as desired. Thus the answer is precisely $\boxed{\text{(D)} \ 23}$.

 


	\item (\textit{AHSME 1985}) Consider the graphs of $y = Ax^2$ and $y^2+3 = x^2+4y$, where $A$ is a positive constant and $x$ and $y$ are real variables. In how many points do the two graphs intersect?


 $\mathrm{(A) \ }\text{exactly }4 \qquad \mathrm{(B) \ }\text{exactly }2 \qquad$


 $\mathrm{(C) \  }\text{at least }1,\text{ but the number varies for different positive values of }A \qquad$


 $\mathrm{(D) \  }0\text{ for at least one positive value of }A \qquad \mathrm{(E) \ }\text{none of these}$


 
 \textbf{Answer: }A



\textbf{Solution 1}

Substituting $y = Ax^2$ into the equation $y^2+3 = x^2+4y$ gives 
 \begin{align*}\left(Ax^2\right)+3 = x^2+4\left(Ax^2\right) &\iff A^2x^4+3 = x^2+4Ax^2 \\ &\iff A^2x^4-\left(4A+1\right)x^2+3 = 0 \\ &\iff x^2 = \frac{4A+1 \pm \sqrt{4A^2+8A+1}}{2A^2} \\ &\text{(using the quadratic formula)}.\end{align*}
 Now observe that since $A$ is positive, $4A^2+8A+1$ is also positive, so the square root will always give two distinct real values. Moreover, 
 \[\left(4A+1\right)^2 = 16A^2+8A+1 > 4A^2+8A+1,\] so $4A+1-\sqrt{4A^2+8A+1} > 0$, meaning that both solutions for $x^2$ are positive. Hence both solutions will give $2$ distinct values of $x$ (the positive and negative square roots), and each of these will correspond to a distinct point of intersection of the graphs, so there are $2 \cdot 2 = \boxed{\text{(A) exactly} \ 4}$ points of intersection.

 



\textbf{Solution 2}

Firstly, note that $y = Ax^2$ is an upward-facing parabola (since $A > 0$) whose vertex is at the origin. We now manipulate the equation of the second graph as follows: 
 \begin{align*}y^2+3 = x^2+4y &\iff y^2-4y+3-x^2 = 0 \\ &\iff y^2-4y+4-x^2 = 1 \\ &\iff \frac{(y-2)^2}{1}-\frac{(x-0)^2}{1} = 1,\end{align*}
 showing that it is a vertical (upward- and downward-opening) hyperbola with center $(0,2)$ and asymptotes $y=x+1$ and $y=-x+1$. It therefore remains to consider graphically where the parabola will intersect the hyperbola.

 On the lower branch of the hyperbola, the maximum point is $(0,2-1) = (0,1)$, which is above the vertex of the parabola. Therefore, by continuity and the symmetry of both the parabola and the hyperbola in the $y$-axis, there are always exactly $2$ intersection points here.

 For the top branch, as it approaches the asymptote $y = x+1$, its slope also approaches that of this asymptote, which is $1$. However, for any upward-opening parabola, the slope approaches infinity as $x$ does, so no matter how small $A$ is (i.e. how 'flat' the parabola is), the parabola will eventually overtake the hyperbola, giving a point of intersection with positive $x$-coordinate. As above, symmetry gives another point of intersection with negative $x$-coordinate, so that there are $2$ intersection points with this branch too.

 Thus there are a total of $2+2 = \boxed{\text{(A) exactly} \ 4}$ intersection points.

 


	\item (\textit{AHSME 1985}) A wooden cube with edge length $n$ units (where $n$ is an integer $>2$) is painted black all over. By slices parallel to its faces, the cube is cut into $n^3$ smaller cubes each of unit edge length. If the number of smaller cubes with just one face painted black is equal to the number of smaller cubes completely free of paint, what is $n$?


 $\mathrm{(A)\ } 5 \qquad \mathrm{(B) \ }6 \qquad \mathrm{(C) \  } 7 \qquad \mathrm{(D) \  } 8 \qquad \mathrm{(E) \  }\text{none of these}$


 
 \textbf{Answer: }D



\textbf{Solution}

Observe that if we remove the outer layer of unit cubes from the entire cube, what remains is a smaller cube of side length $(n-2)$, which contains all of the unpainted cubes and no others. This shows that there are exactly $(n-2)^3$ unpainted cubes. Similarly, taking one face of the cube and removing the outer edge leaves a square of side length $(n-2)$ containing all of the cubes on that face with exactly one face painted. Making the same argument for the other $5$ faces as well, we deduce that there are a total of $6(n-2)^2$ cubes with only one face painted. 

 Accordingly, we require 
 \begin{align*}(n-2)^3 = 6(n-2)^2 &\iff n-2 = 6 \qquad \text{(as } n > 2\text{, so } n-2 \neq 0\text{)} \\ &\iff n = \boxed{\text{(D)} \ 8}.\end{align*}


 


	\item (\textit{AHSME 1985}) How many integers $x$ satisfy the equation 
 \[\left(x^2-x-1\right)^{x+2} = 1?\]


 $\mathrm{(A)\ } 2 \qquad \mathrm{(B) \ }3 \qquad \mathrm{(C) \  } 4 \qquad \mathrm{(D) \  } 5 \qquad \mathrm{(E) \  }\text{none of these}$


 
 \textbf{Answer: }C



\textbf{Solution}

We recall that for real numbers $a$ and $b$, there are exactly $3$ ways in which we can have $a^b = 1$, namely $a = 1$; $b = 0$ and $a \neq 0$; or $a = -1$ and $b$ is an even integer.

 The first case therefore gives 
 \begin{align*}x^2-x-1 = 1 &\iff x^2-x-2 = 0 \\&\iff (x-2)(x+1) = 0 \\&\iff x = 2 \text{ or } x = -1.\end{align*}


 Similarly, the second case gives $x+2 = 0$, i.e. $x = -2$, and this indeed gives $x^2-x-1 = 4+2-1 = 5 \neq 0$, so $x = -2$ is a further valid solution.

 Lastly, for the third case, we have 
 \begin{align*}x^2-x-1 = -1 &\iff x^2-x = 0 \\&\iff x(x-1) = 0 \\&\iff x = 0 \text{ or } x = 1,\end{align*}
 but $x = 1$ would give $x+2 = 3$, which is odd, whereas $x = 0$ gives $x+2 = 2$, which is even. Therefore, this case gives only one further solution, namely $x = 0$.

 Accordingly, the possible values of $x = -2$, $-1$, $0$, or $2$, yielding a total of $\boxed{\text{(C)} \ 4}$ solutions.

 


	\item (\textit{AHSME 1985}) In a circle with center $O$, $AD$ is a diameter, $ABC$ is a chord, $BO = 5$ and $\angle ABO = \ \stackrel{\frown}{CD} \ = 60^{\circ}$. Then the length of $BC$ is


 
                    \begin{figure}[H]
                        centering
                        \includegraphics[width=0.5\textwidth]{Problem/AHSME_I_1985_22_Problem_Figure_1.png}
                    \end{figure}
                    


 $\mathrm{(A)\ } 3 \qquad \mathrm{(B) \ }3+\sqrt{3} \qquad \mathrm{(C) \  } 5-\frac{\sqrt{3}}{2} \qquad \mathrm{(D) \  } 5 \qquad \mathrm{(E) \  }\text{none of the above}$


 
 \textbf{Answer: }D



\textbf{Solution}

Since $\angle CAD$ is an angle inscribed in a $60{^\circ}$ arc, we obtain $\angle CAD =\frac{60^{\circ}}{2} = 30^{\circ}$, so $\triangle ABO$ is a $30^{\circ}$-$60^{\circ}$-$90^{\circ}$ right triangle. This gives $AO = BO\sqrt{3} = 5\sqrt{3}$ and $AB = 2BO = 10$, and now since $AD$ is a diameter, $AD = 2AO = 10\sqrt{3}$. The fact that $AD$ is a diameter also means that $\angle ACD = 90^{\circ}$, so $\triangle ACD$ is again a $30^{\circ}$-$60^{\circ}$-$90^{\circ}$ right triangle, yielding $CD = \frac{1}{2}AD = 5\sqrt{3}$ and $AC = \frac{\sqrt{3}}{2}AD = 15$, which finally gives $BC = AC-AB = 15-10 = \boxed{\text{(D)} \ 5}$.

 


	\item (\textit{AHSME 1985}) If 
 \[x = \frac{-1+i\sqrt{3}}{2} \qquad\text{and}\qquad y = \frac{-1-i\sqrt{3}}{2},\] where $i^2 = -1$, then which of the following is not correct?


 $\mathrm{(A)\ } x^5+y^5 = -1 \qquad \mathrm{(B) \ }x^7+y^7 = -1 \qquad \mathrm{(C) \  } x^9+y^9 = -1 \qquad$


 $\mathrm{(D) \  } x^{11}+y^{11} = -1 \qquad \mathrm{(E) \  }x^{13}+y^{13} = -1$


 
 \textbf{Answer: }C



\textbf{Solution 1}

We can write 
 \begin{align*}&x = \cos\left(\frac{2\pi}{3}\right)+i\sin\left(\frac{2\pi}{3}\right) = e^{\frac{2\pi}{3}i}, \text{ and} \\ &y = \cos\left(-\frac{2\pi}{3}\right)+i\sin\left(-\frac{2\pi}{3}\right) = e^{-\frac{2\pi}{3}i},\end{align*}
 which gives 
 \begin{align*}&x^k = e^{\frac{2\pi k}{3}i} = \cos\left(\frac{2\pi k}{3}\right)+i\sin\left(\frac{2\pi k}{3}\right), \text{ and} \\ &y^k = e^{-\frac{2\pi k}{3}i} = \cos\left(-\frac{2\pi k}{3}\right)+i\sin\left(-\frac{2\pi k}{3}\right) = \cos\left(\frac{2\pi k}{3}\right)-i\sin\left(\frac{2\pi k}{3}\right),\end{align*}
 using the fact that $\cos$ is an even function and $\sin$ is an odd function.

 Accordingly, 
 \[x^k+y^k = 2\cos\left(\frac{2\pi k}{3}\right),\] and upon substituting the values $k = 5,7,9,11,13$ from the answer choices, we find that $x^k+y^k = -1$ for all such values except $k = 9$, where $x^9+y^9 = 2\cos(6\pi) = 2 \neq -1$. Thus the answer is $\boxed{\text{(C)} \ x^9+y^9 = -1}$.

 



\textbf{Solution 2}

Notice that $x+y = -1$ and $xy = 1$, so 
 \begin{align*}1 &=(x+y)^2 \\ &= x^2+2xy+y^2 \\ &=x^2+y^2+2,\end{align*}
 and hence $x^2+y^2=-1$. Similarly, 
 \begin{align*}x^3+y^3 &= (x+y)^3-3x^2y-3xy^2 \\ &= \left(-1\right)^3-3xy(x+y) \\ &= -1-3(1)\left(-1\right) \\ &=2,\end{align*}
 and 
 \begin{align*}x^5+y^5 &= (x+y)^5-5x^4y-5xy^4-10x^3y^2-10x^2y^3 \\ &= \left(-1\right)^5-5xy\left(x^3+y^3+2xy\left(x+y\right)\right) \\ &= -1-5(1)\left(2+2\left(1\right)\left(-1\right)\right) \\ &=-1.\end{align*}


 Now let $z_n = x^{2n+1}+y^{2n+1}$. Then, using the results $x^2+y^2 = -1$ and $xy = 1$ from above, we obtain 
 \begin{align*}-z_n &= \left(x^{2n+1}+y^{2n+1}\right)\left(x^2+y^2\right) \\ &= x^{2n+3}+y^{2n+3}+x^2y^{2n+1}+x^{2n+1}y^2 \\ &= x^{2n+3}+y^{2n+3}+\left(xy\right)^2\left(x^{2n-1}+y^{2n-1}\right) \\ &= x^{2n+3}+y^{2n+3}+x^{2n-1}+y^{2n-1} \\ &= z_{n+1}+z_{n-1}.\end{align*}


 Again from above, $z_1 = 2$ and $z_2 = -1$, so 
 \begin{align*}1 &= -z_2 \\ &= z_3+z_1 \\ &= x^7+y^7+2,\end{align*}
 giving $z_3 = x^7+y^7 = -1$. Similarly, 
 \begin{align*}1 &= -z_3 \\ &= z_4+z_2 \\ &= x^9+y^9-1,\end{align*}
 giving $z_4 = x^9+y^9 = 2 \neq -1$, meaning that the answer must be $\text{(C)}$. To confirm this, we further note that 
 \begin{align*}-2 &= -z_4 \\ &= z_5+z_3 \\ &= x^{11}+y^{11}-1,\end{align*}
 giving $z_5 = x^{11}+y^{11} = -1$, and finally 
 \begin{align*}1 &= -z_5 \\ &= z_6+z_4 \\ &= x^{13}+y^{13}+2,\end{align*}
 giving $x^{13}+y^{13} = -1$, which shows that the only false statement is indeed $\boxed{\text{(C)} \ x^9+y^9 = -1}$.

 


	\item (\textit{AHSME 1985}) A non-zero digit is chosen in such a way that the probability of choosing digit $d$ is $\log_{10}{(d+1)}-\log_{10}{d}$. The probability that the digit $2$ is chosen is exactly $1/2$ the probability that the digit chosen is in the set


 $\mathrm{(A)\ } \{2,3\} \qquad \mathrm{(B) \ }\{3,4\} \qquad \mathrm{(C) \  } \{4,5,6,7,8\} \qquad \mathrm{(D) \  } \{5,6,7,8,9\} \qquad \mathrm{(E) \  }\{4,5,6,7,8,9\}$


 
 \textbf{Answer: }C



\textbf{Solution}

We have $\log_{10}{(d+1)}-\log_{10}{d} = \log_{10}{\left(\frac{d+1}{d}\right)}$, so the probability of choosing $2$ is $\log_{10}{\left(\frac{3}{2}\right)}$. The probability that the digit chosen is in the set must therefore be 
 \begin{align*}2\log_{10}{\left(\frac{3}{2}\right)} = &\log_{10}{\left(\left(\frac{3}{2}\right)^2\right)} \\ = &\log_{10}{\left(\frac{9}{4}\right)} \\ = &\log_{10}{9}-\log_{10}{4} \\ = &\left(\log_{10}{9}-\log_{10}{8}\right)+\left(\log_{10}{8}-\log_{10}{7}\right)+\left(\log_{10}{7}-\log_{10}{6}\right) \\ &+\left(\log_{10}{6}-\log_{10}{5}\right)+\left(\log_{10}{5}-\log_{10}{4}\right),\end{align*}
which, by definition, is the probability that the digit chosen is in the set $\boxed{\text{(C)} \ \{4,5,6,7,8\}}$.

 


	\item (\textit{AHSME 1985}) The volume of a certain rectangular solid is $8$ cm<sup>3</sup>, its total surface area is $32$ cm<sup>2</sup>, and its three dimensions are in geometric progression. The sums of the lengths in cm of all the edges of this solid is


 $\mathrm{(A)\ } 28 \qquad \mathrm{(B) \ }32 \qquad \mathrm{(C) \  } 36 \qquad \mathrm{(D) \  } 40 \qquad \mathrm{(E) \  }44$


 
 \textbf{Answer: }B



\textbf{Solution 1}

As the dimensions are in geometric progression, let them be $\frac{b}{r}$, $b$, and $br$ cm, so the volume is $\left(\frac{b}{r}\right)(b)(br) = b^3$, giving $b^3 = 8$ and thus $b = 2$. The surface area condition now yields 
 \begin{align*}2\left(\frac{2}{r}\right)(2)+2(2)(2r)+2(2r)\left(\frac{2}{r}\right) = 32 &\iff \frac{8}{r}+8+8r = 32 \\ &\iff r+\frac{1}{r} = 3 \\ &\iff r^2-3r+1 = 0 \\ &\iff r = \frac{3 \pm \sqrt{5}}{2}.\end{align*}


 Since 
 \[\frac{3-\sqrt{5}}{2} = \frac{1}{\left(\frac{3+\sqrt{5}}{2}\right)},\] the two possible values of $r$ do not actually give different dimensions, but merely determine whether they are in increasing or decreasing order. Therefore, without loss of generality, we take $r = \frac{3+\sqrt{5}}{2}$, giving the dimensions (in cm) as 
 \begin{align*}&\frac{2}{\left(\frac{3+\sqrt{5}}{2}\right)}, 2, \text{ and } 2\left(\frac{3+\sqrt{5}}{2}\right) \\ &= \frac{4}{3+\sqrt{5}}, 2, \text{ and } 3+\sqrt{5} \\ &= 3-\sqrt{5}, 2, \text{ and } 3+\sqrt{5}.\end{align*}


 As there are $4$ edges with each of these distinct lengths, it follows that the sum of all the edge lengths (in cm) is 
 \begin{align*}4\left(3-\sqrt{5}+2+3+\sqrt{5}\right) &= 4(8) \\ &= \boxed{\text{(B)} \ 32}.\end{align*}


 



\textbf{Solution 2}

Similarly to in Solution 1, we let the dimensions (in cm) be $b$, $br$, and $br^2$, so that the volume condition gives $8 = b^3r^3 = (br)^3$ and thus $br = 2$. The surface area condition now becomes 
 \begin{align*}2(b)(br)+2(br)\left(br^2\right)+2\left(br^2\right)(b) = 32 &\iff b^2r+b^2r^2+b^2r^3 = 16 \\&\iff br\left(b+br+br^2\right) = 16,\end{align*}
 so substituting $br = 2$ from above immediately gives 
 \[b+br+br^2 = \frac{16}{2} = 8,\] and hence, without needing to actually compute the dimensions, we deduce that the sum of the edge lengths (in cm) is 
 \[4b + 4br + 4br^2 = 4(8) = \boxed{\text{(B)} \ 32}.\]

 


	\item (\textit{AHSME 1985}) Find the least positive integer $n$ for which $\frac{n-13}{5n+6}$ is a non-zero reducible fraction.


 $\mathrm{(A)\ } 45 \qquad \mathrm{(B) \ }68 \qquad \mathrm{(C) \  } 155 \qquad \mathrm{(D) \  } 226 \qquad \mathrm{(E) \  }\text{none of these}$


 
 \textbf{Answer: }E



\textbf{Solution}

For the fraction to be reducible, the greatest common factor of the numerator and the denominator must be greater than $1$. Using the \href{https://artofproblemsolving.com/wiki/index.php?title=Euclidean\_algorithm}{Euclidean algorithm}, we compute 
 \begin{align*}\gcd\left(5n+6,n-13\right) &= \gcd\left(5n+6-5(n-13),n-13\right) \\ &= \gcd\left(71,n-13\right).\end{align*}
Since $71$ is prime, it follows that this GCD will be $1$ unless $(n-13)$ is a multiple of $71$, which first occurs when $n = 71+13 = 84$, so the answer is $\boxed{\text{(E) none of these}}$.

 


	\item (\textit{AHSME 1985}) Consider a sequence $x_1,x_2,x_3,\dotsc$ defined by:

 \begin{align*}&x_1 = \sqrt[3]{3}, \\ &x_2 = \left(\sqrt[3]{3}\right)^{\sqrt[3]{3}},\end{align*}

and in general

 \[x_n = \left(x_{n-1}\right)^{\sqrt[3]{3}} \text{ for } n > 1.\]


 What is the smallest value of $n$ for which $x_n$ is an integer?


 $\mathrm{(A)\ } 2 \qquad \mathrm{(B) \ }3 \qquad \mathrm{(C) \  } 4 \qquad \mathrm{(D) \  } 9 \qquad \mathrm{(E) \  }27$


 
 \textbf{Answer: }C



\textbf{Solution 1 (Rigorous Proof)}

Firstly, we will show by induction that 
 \[x_n = \left(\sqrt[3]{3}\right)^{\left(\left(\sqrt[3]{3}\right)^{n-1}\right)}.\] For the base case, we indeed have 
 \begin{align*}x_1 &= \sqrt[3]{3} \\ &= \left(\sqrt[3]{3}\right)^1 \\ &= \left(\sqrt[3]{3}\right)^{\left(\left(\sqrt[3]{3}\right)^0\right)},\end{align*}
 and for the inductive step, if our claim is true for $x_n$, then 
 \begin{align*}x_{n+1} &= \left(x_n\right)^{\sqrt[3]{3}} \\ &= \left(\left(\sqrt[3]{3}\right)^{\left(\left(\sqrt[3]{3}\right)^{n-1}\right)}\right)^{\sqrt[3]{3}} \\ &= \left(\sqrt[3]{3}\right)^{\left(\left(\sqrt[3]{3}\right)^{n-1}\cdot\sqrt[3]{3}\right)} \\ &= \left(\sqrt[3]{3}\right)^{\left(\left(\sqrt[3]{3}\right)^n\right)},\end{align*}
 which completes the proof.

 We now rewrite our formula for $x_n$ as follows:
 \begin{align*}x_n &= \left(\sqrt[3]{3}\right)^{\left(\left(\sqrt[3]{3}\right)^{n-1}\right)} \\ &= \left(\sqrt[3]{3}\right)^{\left(3^{\frac{n-1}{3}}\right)} \\ &=3^{\left(\frac{1}{3} \cdot 3^{\frac{n-1}{3}}\right)} \\ &= 3^{\left(3^{\left(\frac{n-1}{3}-1\right)}\right)} \\ &= 3^{\left(3^{\left(\frac{n-4}{3}\right)}\right)},\end{align*}
and as $3$ is not a perfect power, we deduce that $x_n$ is an integer if and only if the exponent, $3^{\left(\frac{n-4}{3}\right)}$, is itself an integer. By precisely the same argument, this reduces to $\frac{n-4}{3}$ being an integer, so the smallest possible (positive) value of $n$ is $\boxed{\text{(C)} \ 4}$.

 



\textbf{Solution 2 (Hit and Trial)}

Observe that $x_2 = \sqrt[3]{3}^{\sqrt[3]{3}}$ is not an integer and $x_3 = \sqrt[3]{3}^{\sqrt[3]{9}}$ is also not an integer. However, $x_4 = \sqrt[3]{3}^{\sqrt[3]{27}} = \sqrt[3]{3}^3 = 3$ is an integer. Hence at $n = 4$ we hit an integer for the first time, and thus our answer is $\boxed{\text{(C)} \ 4}$.

 


	\item (\textit{AHSME 1985}) In $\triangle ABC$, we have $\angle C = 3\angle A$, $a = 27$ and $c = 48$. What is $b$?


 
                    \begin{figure}[H]
                        centering
                        \includegraphics[width=0.5\textwidth]{Problem/AHSME_I_1985_28_Problem_Figure_1.png}
                    \end{figure}
                    


 $\mathrm{(A)\ } 33 \qquad \mathrm{(B) \ }35 \qquad \mathrm{(C) \  } 37 \qquad \mathrm{(D) \  } 39 \qquad \mathrm{(E) \  }\text{not uniquely determined}$


 
 \textbf{Answer: }B



\textbf{Solution 1}

Let $\angle A = x^{\circ}$, so $\angle C = 3x^{\circ}$, and thus $\angle B = \left(180-4x\right)^{\circ}$. Now let $D$ be a point on side $AB$ such that $\angle ACD = x^{\circ}$, so $\angle BCD = 3x^{\circ}-x^{\circ} = 2x^{\circ}$, which gives 
 \[\angle CDB = 180^{\circ}-2x^{\circ}-\left(180-4x\right)^{\circ} = 2x^{\circ},\] meaning that $\triangle CDB$ and $\triangle CDA$ are both isosceles, with $BC = BD$ and $AD = CD$. In particular, $BD = BC = 27$ and $CD = AD = AB-BD = 48-27 = 21$. Hence by \href{https://artofproblemsolving.com/wiki/index.php?title=Stewart\%27s\_theorem}{Stewart's theorem} on triangle $ABC$, 
 \begin{align*}&BD \cdot AB \cdot AD + CD^2 \cdot AB = AC^2 \cdot BD + BC^2 \cdot AD \\ &\iff 27 \cdot 48 \cdot 21 + 21^2 \cdot 48 = AC^2 \cdot 27 + 27^2 \cdot 21 \\ &\iff AC^2 = \frac{27(21)(48-27) + 21^2 \cdot 48}{27} \\ &\iff AC^2 = \frac{21^2(27+48)}{27} \\ &\iff AC^2 = \frac{21^2 \cdot 25}{9} \\ &\iff AC = \frac{21 \cdot 5}{3} \qquad \text{(as } AC > 0\text{)} \\ &\iff AC = \boxed{\text{(B)} \ 35}.\end{align*}


 



\textbf{Solution 2}

We apply the law of sines in the form 
 \[\frac{\sin(A)}{a} = \frac{\sin(C)}{c},\] yielding 
 \[\frac{\sin(A)}{27} = \frac{\sin(3A)}{48} \iff 9\sin(3A) = 16\sin(A).\]

 Now, the angle sum and double angle identities give 
 \begin{align*}\sin(3A) &= \sin(2A+A) \\ &= \sin(2A)\cos(A)+\cos(2A)\sin(A) \\ &= \left(2\sin(A)\cos(A)\right)\cos(A)+\left(\cos^2(A)-\sin^2(A)\right)\sin(A) \\ &= 2\sin(A)\cos^2(A)+\sin(A)\cos^2(A)-\sin^3(A) \\ &= 3\sin(A)\left(1-\sin^2(A)\right)-\sin^3(A) \\ &\text{(using the further identity } \cos^2(\theta)+\sin^2(\theta) = 1\text{)} \\ &= 3\sin(A)-4\sin^3(A).\end{align*}


 Thus our equation becomes 
 \begin{align*}9\left(3\sin(A)-4\sin^3(A)\right) = 16\sin(A) &\iff 27\sin(A)-36\sin^3(A) = 16\sin(A) \\ &\iff 36\sin^3(A) = 11\sin(A) \\ &\iff \sin(A) = 0 \text{ or } \pm\frac{\sqrt{11}}{6}.\end{align*}
Notice, however, that we must have $0^{\circ} < A < 45^{\circ}$, the latter because otherwise $A+3A \geq 180^{\circ}$, which would contradict the fact that $A$ and $3A$ are angles in a (non-degenerate) triangle. This means $\sin(A) > 0$, so the only valid solution is 
 \[\sin(A) = \frac{\sqrt{11}}{6},\] and the fact that $A$ is acute also means $\cos(A) > 0$, so we deduce 
 \begin{align*}\cos(A) &= \sqrt{1-\left(\frac{\sqrt{11}}{6}\right)^2} \\ &= \sqrt{1-\frac{11}{36}} \\ &=\frac{5}{6}.\end{align*}
Accordingly, using the double angle identities again, 
 \begin{align*}\sin(4A) &= \sin(2 \cdot 2A) \\ &= 2\sin(2A)\cos(2A) \\ &= 2\left(2\sin(A)\cos(A)\right)\left(\cos^2(A)-\sin^2(A)\right) \\ & =2\left(2 \cdot \frac{\sqrt{11}}{6} \cdot\frac{5}{6}\right)\left(\left(\frac{5}{6}\right)^2-\left(\frac{\sqrt{11}}{6}\right)^2\right) \\ &= \frac{5\sqrt{11}}{9} \cdot \frac{25-11}{36} \\ &= \frac{5\sqrt{11}}{9} \cdot \frac{7}{18} \\ &= \frac{35\sqrt{11}}{162}.\end{align*}
Finally, the law of sines now gives 
 \begin{align*}\frac{\sin(A)}{27} &= \frac{\sin(B)}{b} \\ &= \frac{\sin(180^{\circ}-3A-A)}{b} \\ &= \frac{\sin(4A)}{b} \qquad \text{(using the identity } \sin\left(180^{\circ}-\theta\right) = \sin(\theta)\text{)},\end{align*}
 so, substituting the above results, 
 \[\frac{\left(\frac{\sqrt{11}}{6}\right)}{27} = \frac{\left(\frac{35\sqrt{11}}{162}\right)}{b} \iff b = \frac{6 \cdot 27 \cdot 35}{162} = \boxed{\text{(B)} \ 35}.\]

 


	\item (\textit{AHSME 1985}) In their base $10$ representations, the integer $a$ consists of a sequence of $1985$ eights and the integer $b$ consists of a sequence of $1985$ fives. What is the sum of the digits of the base $10$ representation of the integer $9ab$?


 $\mathrm{(A)\ } 15880 \qquad \mathrm{(B) \ }17856 \qquad \mathrm{(C) \  } 17865 \qquad \mathrm{(D) \  } 17874 \qquad \mathrm{(E) \  }19851$


 
 \textbf{Answer: }C



\textbf{Solution 1}

By the formula for the sum of a geometric series, 
 \begin{align*}a &= 8 \cdot 10^0 + 8 \cdot 10^1 + \dotsb + 8 \cdot 10^{1984} \\ &= \frac{8\left(10^{1985}-1\right)}{10-1} \\ &= \frac{8\left(10^{1985}-1\right)}{9},\end{align*}
 and similarly 
 \[b = \frac{5\left(10^{1985}-1\right)}{9},\] so 
 \begin{align*}9ab &= 9\cdot\frac{8\left(10^{1985}-1\right)}{9}\cdot\frac{5\left(10^{1985}-1\right)}{9} \\ &= \frac{40\left(10^{1985}-1\right)^2}{9} \\ &= \frac{40\left(10^{3970}-2 \cdot 10^{1985}+1\right)}{9} \\ &= \frac{4\left(10^{3971}-2 \cdot 10^{1986}+10\right)}{9}.\end{align*}


 We now compute the decimal expansion of this expression. Firstly, $10^{3971} = 100 \dotsb 0$, with $1$ one and $3971$ zeroes, and $2 \cdot 10^{1986} = 200 \dotsb 0$, with $1$ two and $1986$ zeroes. Subtracting therefore gives 
 \[10^{3971}-2 \cdot 10^{1986} = 99 \dotsb 9800 \dotsb 0,\] where there are $3971-1986-1 = 1984$ nines followed by $1$ eight and then $1986$ zeroes. Adding $10$ transforms this to $99 \dotsb 9800 \dotsb 010$, now with $1984$ nines followed by $1$ eight, $1984$ zeroes, $1$ one, and a final zero.

 Using long division, and noting that $80 = 8 \cdot 9 + 8$ and $81 = 9 \cdot 9$, it follows that 
 \[\frac{4\left(10^{3971}-2 \cdot 10^{1986}+10\right)}{9} = 11 \dotsb 1088 \dotsb 890,\] with $1984$ ones, $1$ zero, then $1984$ eights, $1$ nine, and a final zero. Lastly, using long multiplication and noting that $9 \cdot 4 = 36$, $8 \cdot 4 = 32$, and $8 \cdot 4 + 3 = 35$, we obtain 
 \[\frac{4\left(10^{3971}-2 \cdot 10^{1986}+10\right)}{9} = 44 \dotsb 4355 \dotsb 560,\] where there are $1984$ fours, $1$ three, $1984$ fives, $1$ six, and a final zero, so the sum of the digits is 
 \begin{align*}1984 \cdot 4 + 3 + 1984 \cdot 5 + 6 + 0 &= 1984 \cdot 9 + 9 \\ &= 1985 \cdot 9 \\ &= \boxed{\text{(C)} \ 17865}.\end{align*}


 



\textbf{Solution 2}

Factoring out the 8 and the 5, we get 
 \[9ab = 10\cdot4\cdot9\cdot111 \dotsb 1 \cdot 111\dotsb1,\] where there are $1985$ ones in both numbers.

 


 This is equal to 
 \[10 \cdot 4 \cdot (1000\dotsb 0 - 1) \cdot 111\dotsb 1\] where there is an initial $1$ followed by $1985$ zeros and the second number has $1985$ ones.

 


 Multiplying then subtracting gives us 
 \[4 \cdot 10 \cdot 111\dotsb 10888 \dotsb 89,\] where there are $1984$ ones and $1984$ eights.

 


 Multiplying one last time, then adding the $0$ at the end gives us 
 \[44 \dotsb 4355 \dotsb 560,\] where there are $1984$ fours, $1$ three, $1984$ fives, $1$ six, and a final zero, so the sum of the digits is 
 \[1984(4+5)+6+3 = 1985 \cdot 9 = \boxed{ \text{ C)     }17865}\]

 


 


	\item (\textit{AHSME 1985}) Let $\left\lfloor x\right\rfloor$ be the greatest integer less than or equal to $x$. Then the number of real solutions to $4x^2-40\left\lfloor x\right\rfloor+51 = 0$ is


 $\mathrm{(A)\ } 0 \qquad \mathrm{(B) \ }1 \qquad \mathrm{(C) \  } 2 \qquad \mathrm{(D) \  } 3 \qquad \mathrm{(E) \  }4$


 
 \textbf{Answer: }E



\textbf{Solution 1}

We rearrange the equation as $4x^2 = 40\left\lfloor x\right\rfloor-51$, where the right-hand side is now clearly an integer, meaning that $4x^2 = n$ for some non-negative integer $n$. Therefore, in the case where $x \geq 0$, substituting $x = \frac{\sqrt{n}}{2}$ gives
 \[40\left\lfloor\frac{\sqrt{n}}{2}\right\rfloor-51 = n.\] To proceed, let $a$ be the unique non-negative integer such that $a \leq \frac{\sqrt{n}}{2} < a+1$, so that 
 \begin{align*}&\left\lfloor \frac{\sqrt{n}}{2}\right\rfloor = a, \text{ and} \\ &4a^2 \leq n < 4a^2+8a+4,\end{align*}
 and our equation reduces to 
 \[40a-51 = n.\]

 The above inequalities therefore become 
 \[4a^2 \leq 40a-51 < 4a^2+8a+4  \iff 4a^2-40a+51 < 0 \text{ and } 4a^2-32a+55 > 0,\] where the first inequality can now be rewritten as $(2a-10)^2 \leq 49$, i.e. $\left\lvert 2a-10\right\rvert \leq 7$. Since $(2a-10)$ is even for all integers $a$, we must in fact have 
 \begin{align*}\left\lvert 2a-10\right\rvert \leq 6 &\iff \left\lvert a-5\right\rvert \leq 3 \\ &\iff 2 \leq a \leq 8.\end{align*}
 The second inequality similarly simplifies to $(2a-8)^2 > 9$, i.e. $\left\lvert 2a-8\right\rvert > 3$. As $(2a-8)$ is even, this is equivalent to 
 \begin{align*}\left\lvert 2a-8 \right\rvert \geq 4 &\iff \left\lvert a-4\right\rvert \geq 2 \\ &\iff a \geq 6 \text{ or } a \leq 2,\end{align*}
 so the values of $a$ satisfying both inequalities are $2$, $6$, $7$, and $8$. Since $n = 40a-51$, each of these distinct values of $a$ gives a distinct solution for $n$, and thus for $x = \frac{\sqrt{n}}{2}$, giving a total of $4$ solutions in the $x \geq 0$ case.

 
As $4$ is already the largest of the answer choices, this suffices to show that the answer is $\text{(E)}$, but for completeness, we will show that the $x < 0$ case indeed gives no other solutions. If $x = -\frac{\sqrt{n}}{2}$ (and so $n > 0$), we require 
 \[40\left\lfloor -\frac{\sqrt{n}}{2}\right\rfloor-51 = n,\] and recalling that $\left\lfloor -x\right\rfloor = -\left\lceil x\right\rceil$ for all $x$, this equation can be rewritten as 
 \[-40\left\lceil \frac{\sqrt{n}}{2}\right\rceil-51 = n.\] Since $n$ is positive, the least possible value of $\left\lceil \frac{\sqrt{n}}{2}\right\rceil$ is $1$, but this means 
 \begin{align*}n &= -40\left\lceil\frac{\sqrt{n}}{2}\right\rceil-51 \\ &\leq -40 \cdot 1 - 51 \\ &= -91,\end{align*}
 which is a contradiction. Therefore the $x < 0$ case indeed gives no further solutions, confirming that the total number of solutions is precisely $\boxed{\text{(E)} \ 4}$.

 



\textbf{Solution 2}

This equation is equivalent to $4x^2 - 40x + 40x - 40\lfloor x \rfloor + 51 = 0$, or $4x^2 - 40x + 51 + 40(x - \lfloor x \rfloor) = 0$.  Let $\{x\} = x - \lfloor x \rfloor$ be the fractional part of $x$.  Then this equation becomes $4x^2 - 40x + 51 + 40\{x\} = 0$.

 Note that for all $x$, $0 \leq \{x\} < 1$.  Therefore, $4x^2 - 40x + 51 \leq 4x^2 - 40x + 51 + 40\{x\} \leq 4x^2 - 40x + 91$, that is $4x^2 - 40x + 51 \leq 0 < 4x^2 - 40x + 91$.  

 First, since $4x^2 - 40x + 51 = (2x - 17)(2x - 3) \leq 0$, it follows that $\frac{3}{2} \leq x \leq \frac{17}{2}$.  

 Also, since $4x^2 - 40x + 91 = (2x - 13)(2x - 7) > 0$, either $x < \frac{7}{2}$ or $x > \frac{13}{2}$.

 Thus either $\frac{3}{2} \leq x < \frac{7}{2}$ or $\frac{13}{2} < x \leq \frac{17}{2}$, so $\lfloor x \rfloor \in \{1, 2, 3, 6, 7, 8\}$.

 Isolating $x^2$ in the original equation yields $x^2 = \frac{40\lfloor x \rfloor - 51}{4}$.  

 If $\lfloor x \rfloor = 1$, then $x^2 = -\frac{11}{4}$.  This does not yield any valid solutions as $x^2$ may not be negative.

 If $\lfloor x \rfloor = 2$, then $x^2 = \frac{29}{4}$ and $x = \frac{\sqrt{29}}{2}$.  This is a valid solution since $\left\lfloor\frac{\sqrt{29}}{2}\right\rfloor = 2$.

 If $\lfloor x \rfloor = 3$, then $x^2 = \frac{69}{4}$ and $x = \frac{\sqrt{69}}{2}$.  This is not a valid solution since $\left\lfloor\frac{\sqrt{69}}{2}\right\rfloor = 4$.

 If $\lfloor x \rfloor = 6$, then $x^2 = \frac{189}{4}$ and $x = \frac{\sqrt{189}}{2}$.  This is a valid solution since $\left\lfloor\frac{\sqrt{189}}{2}\right\rfloor = 6$.

 If $\lfloor x \rfloor = 7$, then $x^2 = \frac{229}{4}$ and $x = \frac{\sqrt{229}}{2}$.  This is a valid solution since $\left\lfloor\frac{\sqrt{229}}{2}\right\rfloor = 7$.

 If $\lfloor x \rfloor = 8$, then $x^2 = \frac{269}{4}$ and $x = \frac{\sqrt{269}}{2}$.  This is a valid solution since $\left\lfloor\frac{\sqrt{269}}{2}\right\rfloor = 8$.

 Therefore, this equation as a total of $\boxed{(\textbf{E})\ 4}$ solutions.

 -j314andrews

 


\end{enumerate}